\documentclass[12pt,a4paper]{article}
\usepackage{mathwizards-ingles}
\title{%
  \vspace{-1cm}
  {\Huge\bfseries\textcolor{mes1}{Mes 1 · Semanas 1--2}}\\[0.3cm]
  {\Large Survival Phrases, Saludos y Vocabulario TPR}\\[0.2cm]
  {\large\color{grissuave} Nivel A0 · El Periodo Silencioso}\\[0.5cm]
  \rule{\textwidth}{1.5pt}
}
\author{}
\date{}

\begin{document}
\maketitle
\thispagestyle{empty}

\begin{cajanivel}
\textbf{Objetivo:} Construir una base fonética sólida y un vocabulario pasivo inicial. En esta etapa \textbf{no se espera producción oral}; el estudiante responde con gestos, reacciones y señalamiento (TPR). El foco está en \textbf{escuchar y reconocer}.
\end{cajanivel}

% ═══════════════════════════════════════════════════════════════════════════════
\section{Chunks: Saludos y Presentaciones}
% ═══════════════════════════════════════════════════════════════════════════════

\begin{cajachunk}[Greetings — Saludos]
\begin{tabularx}{\textwidth}{@{} X l @{}}
\eng{Hello!} / \eng{Hi!} & \esp{¡Hola!} \\
\eng{Good morning.} & \esp{Buenos días.} \\
\eng{Good afternoon.} & \esp{Buenas tardes.} \\
\eng{Good evening.} & \esp{Buenas noches (al llegar).} \\
\eng{Good night.} & \esp{Buenas noches (al despedirse).} \\
\eng{How are you?} & \esp{¿Cómo estás?} \\
\eng{I'm fine, thank you.} & \esp{Estoy bien, gracias.} \\
\eng{I'm good. And you?} & \esp{Estoy bien. ¿Y tú?} \\
\eng{Nice to meet you.} & \esp{Mucho gusto.} \\
\eng{See you later!} & \esp{¡Hasta luego!} \\
\eng{Goodbye! / Bye!} & \esp{¡Adiós!} \\
\end{tabularx}
\end{cajachunk}

\begin{cajachunk}[Self-Introduction — Presentación Personal]
\begin{tabularx}{\textwidth}{@{} X l @{}}
\eng{My name is \underline{\hspace{2cm}}.} & \esp{Mi nombre es \underline{\hspace{2cm}}.} \\
\eng{I am from \underline{\hspace{2cm}}.} & \esp{Soy de \underline{\hspace{2cm}}.} \\
\eng{I am \underline{\hspace{1cm}} years old.} & \esp{Tengo \underline{\hspace{1cm}} años.} \\
\eng{I live in \underline{\hspace{2cm}}.} & \esp{Vivo en \underline{\hspace{2cm}}.} \\
\eng{I speak Spanish.} & \esp{Hablo español.} \\
\eng{I am a student.} & \esp{Soy estudiante.} \\
\eng{I like \underline{\hspace{2cm}}.} & \esp{Me gusta \underline{\hspace{2cm}}.} \\
\end{tabularx}
\end{cajachunk}

% ═══════════════════════════════════════════════════════════════════════════════
\section{Chunks: Vocabulario TPR (Respuesta Física Total)}
% ═══════════════════════════════════════════════════════════════════════════════

\begin{cajachunk}[Classroom Commands — Comandos del Aula]
Estas frases las usa el \textbf{profesor} para dar instrucciones. El estudiante responde \textbf{con el cuerpo}, no con palabras.

\medskip
\begin{multicols}{2}
\begin{tabularx}{\linewidth}{@{} X @{}}
\eng{Stand up.} \esp{Levántate.} \\
\eng{Sit down.} \esp{Siéntate.} \\
\eng{Turn around.} \esp{Date la vuelta.} \\
\eng{Jump!} \esp{¡Salta!} \\
\eng{Walk to the door.} \esp{Camina a la puerta.} \\
\eng{Run!} \esp{¡Corre!} \\
\eng{Stop!} \esp{¡Para!} \\
\eng{Clap your hands.} \esp{Aplaude.} \\
\end{tabularx}

\columnbreak

\begin{tabularx}{\linewidth}{@{} X @{}}
\eng{Point to \underline{\hspace{1.5cm}}.} \esp{Señala \underline{\hspace{1.5cm}}.} \\
\eng{Touch your nose.} \esp{Toca tu nariz.} \\
\eng{Open the book.} \esp{Abre el libro.} \\
\eng{Close the door.} \esp{Cierra la puerta.} \\
\eng{Pick up the pen.} \esp{Recoge el bolígrafo.} \\
\eng{Put it down.} \esp{Ponlo abajo.} \\
\eng{Look at me.} \esp{Mírame.} \\
\eng{Listen!} \esp{¡Escucha!} \\
\end{tabularx}
\end{multicols}
\end{cajachunk}

\begin{cajachunk}[Action Verbs — Verbos de Acción (Banco TPR)]
\begin{multicols}{3}
\begin{itemize}[nosep, leftmargin=0.6cm]
  \item \eng{run} \esp{correr}
  \item \eng{walk} \esp{caminar}
  \item \eng{jump} \esp{saltar}
  \item \eng{sit} \esp{sentarse}
  \item \eng{stand} \esp{pararse}
  \item \eng{stop} \esp{parar}
  \item \eng{go} \esp{ir}
  \item \eng{come} \esp{venir}
  \item \eng{open} \esp{abrir}
  \item \eng{close} \esp{cerrar}
  \item \eng{touch} \esp{tocar}
  \item \eng{point} \esp{señalar}
  \item \eng{look} \esp{mirar}
  \item \eng{listen} \esp{escuchar}
  \item \eng{eat} \esp{comer}
  \item \eng{drink} \esp{beber}
  \item \eng{read} \esp{leer}
  \item \eng{write} \esp{escribir}
  \item \eng{give} \esp{dar}
  \item \eng{take} \esp{tomar}
  \item \eng{put} \esp{poner}
  \item \eng{show} \esp{mostrar}
  \item \eng{say} \esp{decir}
  \item \eng{help} \esp{ayudar}
\end{itemize}
\end{multicols}
\end{cajachunk}

% ═══════════════════════════════════════════════════════════════════════════════
\section{Vocabulario Base: Números, Días y Colores}
% ═══════════════════════════════════════════════════════════════════════════════

\begin{cajachunk}[Numbers 1--20 — Números]
\begin{multicols}{4}
\begin{enumerate}[nosep, leftmargin=1cm]
  \item \eng{one}
  \item \eng{two}
  \item \eng{three}
  \item \eng{four}
  \item \eng{five}
  \item \eng{six}
  \item \eng{seven}
  \item \eng{eight}
  \item \eng{nine}
  \item \eng{ten}
  \item \eng{eleven}
  \item \eng{twelve}
  \item \eng{thirteen}
  \item \eng{fourteen}
  \item \eng{fifteen}
  \item \eng{sixteen}
  \item \eng{seventeen}
  \item \eng{eighteen}
  \item \eng{nineteen}
  \item \eng{twenty}
\end{enumerate}
\end{multicols}
\end{cajachunk}

\begin{multicols}{2}

\begin{cajachunk}[Days of the Week — Días]
\begin{tabularx}{\linewidth}{@{} X l @{}}
\eng{Monday} & \esp{lunes} \\
\eng{Tuesday} & \esp{martes} \\
\eng{Wednesday} & \esp{miércoles} \\
\eng{Thursday} & \esp{jueves} \\
\eng{Friday} & \esp{viernes} \\
\eng{Saturday} & \esp{sábado} \\
\eng{Sunday} & \esp{domingo} \\
\end{tabularx}
\end{cajachunk}

\columnbreak

\begin{cajachunk}[Colors — Colores]
\begin{tabularx}{\linewidth}{@{} X l @{}}
\eng{red} & \esp{rojo} \\
\eng{blue} & \esp{azul} \\
\eng{green} & \esp{verde} \\
\eng{yellow} & \esp{amarillo} \\
\eng{black} & \esp{negro} \\
\eng{white} & \esp{blanco} \\
\eng{orange} & \esp{naranja} \\
\eng{purple} & \esp{morado} \\
\eng{pink} & \esp{rosado} \\
\eng{brown} & \esp{marrón} \\
\end{tabularx}
\end{cajachunk}

\end{multicols}

% ═══════════════════════════════════════════════════════════════════════════════
\section{Chunks: Frases de Supervivencia en el Aula Virtual}
% ═══════════════════════════════════════════════════════════════════════════════

\begin{cajachunk}[Survival Phrases — Para Usar en las Sesiones]
\begin{tabularx}{\textwidth}{@{} X l @{}}
\eng{I don't understand.} & \esp{No entiendo.} \\
\eng{Can you repeat, please?} & \esp{¿Puedes repetir, por favor?} \\
\eng{What does \underline{\hspace{2cm}} mean?} & \esp{¿Qué significa \underline{\hspace{2cm}}?} \\
\eng{How do you say \underline{\hspace{2cm}}?} & \esp{¿Cómo se dice \underline{\hspace{2cm}}?} \\
\eng{Slowly, please.} & \esp{Despacio, por favor.} \\
\eng{One more time, please.} & \esp{Una vez más, por favor.} \\
\eng{Yes. / No.} & \esp{Sí. / No.} \\
\eng{Thank you. / Thanks!} & \esp{Gracias.} \\
\eng{Sorry.} / \eng{Excuse me.} & \esp{Perdón. / Disculpa.} \\
\eng{I have a question.} & \esp{Tengo una pregunta.} \\
\end{tabularx}
\end{cajachunk}

% ═══════════════════════════════════════════════════════════════════════════════
\section{Gramática Emergente: El Verbo \textit{To Be}}
% ═══════════════════════════════════════════════════════════════════════════════

\begin{cajagrammar}[El verbo \textit{to be} — Presente Simple]
En esta etapa, la gramática se adquiere \textbf{inconscientemente} a través del input. Esta sección es solo una \textbf{referencia} para que el estudiante reconozca el patrón cuando lo encuentre.

\medskip
\begin{tabularx}{\textwidth}{@{} l l X @{}}
\toprule
\textbf{Sujeto} & \textbf{Forma} & \textbf{Ejemplo} \\
\midrule
I & \eng{am} & \textit{I \textbf{am} a student.} \\
You & \eng{are} & \textit{You \textbf{are} my friend.} \\
He / She / It & \eng{is} & \textit{She \textbf{is} a teacher.} \\
We & \eng{are} & \textit{We \textbf{are} happy.} \\
They & \eng{are} & \textit{They \textbf{are} in the classroom.} \\
\bottomrule
\end{tabularx}

\medskip
\textbf{Contracciones} (las formas que realmente escucharás):
\begin{multicols}{3}
\begin{itemize}[nosep, leftmargin=0.6cm]
  \item \eng{I'm} = I am
  \item \eng{You're} = You are
  \item \eng{He's} = He is
  \item \eng{She's} = She is
  \item \eng{It's} = It is
  \item \eng{We're} = We are
  \item \eng{They're} = They are
\end{itemize}
\end{multicols}
\end{cajagrammar}

% ═══════════════════════════════════════════════════════════════════════════════
\section{Gramática Emergente: \textit{This / That / These / Those}}
% ═══════════════════════════════════════════════════════════════════════════════

\begin{cajagrammar}[Demostrativos — Señalar Objetos (clave para TPR)]
\begin{tabularx}{\textwidth}{@{} l l X @{}}
\toprule
 & \textbf{Singular} & \textbf{Plural} \\
\midrule
Cerca & \eng{this} \esp{este/esta} & \eng{these} \esp{estos/estas} \\
Lejos & \eng{that} \esp{ese/aquel} & \eng{those} \esp{esos/aquellos} \\
\bottomrule
\end{tabularx}

\medskip
\begin{cajaejemplo}[Ejemplos en Contexto]
\begin{itemize}[nosep, leftmargin=0.8cm]
  \item \textit{\textbf{This} is a pen.} \esp{Esto es un bolígrafo.}
  \item \textit{\textbf{That} is the door.} \esp{Eso es la puerta.}
  \item \textit{\textbf{These} are my books.} \esp{Estos son mis libros.}
  \item \textit{\textbf{Those} are your chairs.} \esp{Esas son tus sillas.}
\end{itemize}
\end{cajaejemplo}
\end{cajagrammar}

% ═══════════════════════════════════════════════════════════════════════════════
\section{Oraciones Listas para Minar (Modelo de Tarjeta Anki)}
% ═══════════════════════════════════════════════════════════════════════════════

\begin{cajamineria}[Oraciones \textit{i+1} — Una sola incógnita por oración]
Estas oraciones están diseñadas para que un estudiante A0 tenga \textbf{un solo elemento nuevo} que aprender. La palabra objetivo está \underline{subrayada}.

\medskip
\begin{enumerate}[leftmargin=1.2cm]
  \item \textit{My name is Ana. \underline{Nice to meet you}.}\\
    {\small\textbf{Objetivo:} la expresión ``Nice to meet you'' como unidad.}

  \item \textit{Stand up and \underline{clap} your hands.}\\
    {\small\textbf{Objetivo:} el verbo ``clap'' (aplaudir).}

  \item \textit{I am from Colombia. I \underline{speak} Spanish.}\\
    {\small\textbf{Objetivo:} el verbo ``speak'' (hablar).}

  \item \textit{The book is \underline{on} the table.}\\
    {\small\textbf{Objetivo:} la preposición ``on'' (sobre).}

  \item \textit{Point to the \underline{window}, please.}\\
    {\small\textbf{Objetivo:} ``window'' (ventana).}

  \item \textit{What color is this? It's \underline{purple}.}\\
    {\small\textbf{Objetivo:} el color ``purple'' (morado).}

  \item \textit{Today is \underline{Wednesday}. Tomorrow is Thursday.}\\
    {\small\textbf{Objetivo:} ``Wednesday'' (miércoles).}

  \item \textit{She is a \underline{teacher}. He is a student.}\\
    {\small\textbf{Objetivo:} ``teacher'' (profesora).}

  \item \textit{I don't understand. Can you repeat, \underline{please}?}\\
    {\small\textbf{Objetivo:} la palabra ``please'' como unidad de cortesía.}

  \item \textit{Look at me and \underline{listen}.}\\
    {\small\textbf{Objetivo:} el verbo ``listen'' (escuchar).}
\end{enumerate}
\end{cajamineria}

\vfill
\begin{center}
\rule{0.5\textwidth}{0.5pt}\\[0.3cm]
{\small\color{grissuave} Curso de Inmersión en Inglés · Mes 1 · Semanas 1--2 · Nivel A0}
\end{center}

\end{document}
